\documentclass[11pt]{article}
\usepackage[paperwidth=220mm,paperheight=320mm,left=16mm,right=16mm,top=18mm,bottom=18mm,headheight=16pt,headsep=10pt,footskip=14pt]{geometry}
\usepackage{fix-cm}
\usepackage{fontspec,amsmath,amssymb,graphicx,array,multirow,tikz}
\usetikzlibrary{decorations.pathreplacing}
\usepackage[unicode,hidelinks]{hyperref}
\usepackage{polyglossia}
\setmainlanguage{latin}
\setotherlanguages{arabic,greek,hebrew,syriac}
\setmainfont{Linux Libertine O}[Ligatures=TeX]
\setsansfont{Linux Biolinum O}[Ligatures=TeX]
\newfontfamily\greekfont{FreeSerif}[Script=Greek]
\newfontfamily\arabicfont{Amiri}[Script=Arabic]
\newfontfamily\hebrewfont{Frank Ruhl Hofshi}[Script=Hebrew]
\newfontfamily\syriacfont{Segoe UI Historic}[Script=Syriac]
\newfontfamily\CapFont{Noto Serif}
\newfontfamily\MoonFont{FreeSerif}
\newcommand{\MoonSym}{\text{\MoonFont ☾}}\newcommand{\SunSym}{\text{\MoonFont ⊙}}
\newfontfamily\EthiopicFont{Ebrima}
\newfontfamily\CopticFont{FreeSerif}
\newcommand{\textcoptic}[1]{{\CopticFont\bfseries #1}}
\newcommand{\textethiopic}[1]{{\EthiopicFont #1}}
\newfontfamily\NallinoSigns{NallinoSigns.otf}[Path=./fonts/]
\newcommand{\AbjadZero}{{\NallinoSigns\symbol{"E000}}}
% the sign of right ascension: the left half of the A (its thin stroke and its bar) joined to the stem of the R
\usepackage{trimclip}
\newcommand{\ARsign}{\mbox{\clipbox{0pt 0pt {.48\width} 0pt}{A}\kern-.10em R}}
\hypersetup{pdftitle={Al-Battani, Opus astronomicum. Nallino, pars II: frontispicium, praefatio, conspectus, addenda. S05},pdfauthor={Al-Battani; Carlo Alfonso Nallino (Latin translation and notes)}}
\makeatletter
\def\ps@sourceedition{%
\def\@oddhead{\vbox{\hbox to\textwidth{\fontsize{8}{10}\selectfont C. A. NALLINO — PARS SECUNDA\hfil\leftmark}\vskip3pt\hrule height.25pt}}%
\let\@evenhead\@oddhead
\def\@oddfoot{\hbox to\textwidth{\fontsize{7}{9}\selectfont\rightmark\hfil\thepage}}%
\let\@evenfoot\@oddfoot}
\makeatother
\pagestyle{sourceedition}
\setlength{\parindent}{0pt}\setlength{\parskip}{0pt}
\newcommand{\BeginSource}[2]{\clearpage\markboth{PAG. #2}{#1}\hypertarget{#1}{}\pdfbookmark[0]{#1 — p. #2}{bk-#1}\fontsize{11.1}{13.8}\selectfont}
\tracinglostchars=2
\newcommand{\Name}[1]{{\addfontfeatures{LetterSpace=4.0}#1}}
\newcommand{\Indent}{\hspace*{1.6em}}
% an Arabic run inside a Latin line: neither its descenders nor its vowel signs above (shadda, damma) open the line
% (the printed lines are evenly spaced)
\newcommand{\ArabicRun}[1]{\smash{\textarabic{#1}}}
\newcommand{\qfrac}[2]{{}^{#1}\!/_{{#2}}}
\newcommand{\sa}{\textsuperscript{\textit{a}}}\newcommand{\sm}{\textsuperscript{\textit{m}}}\newcommand{\sd}{\textsuperscript{\textit{d}}}
\newcommand{\RawBlock}[1]{\par\vspace{3pt}\noindent#1\par\vspace{3pt}}
\newcommand{\CenterLine}[1]{\makebox[\linewidth][c]{#1}}
\newsavebox{\LBox}\newcommand{\LFit}[3]{\sbox{\LBox}{#3}\resizebox{#1}{#2}{\usebox{\LBox}}}
\newlength{\FullSourceWidth}\newlength{\GutterOffset}\newif\ifNumbersLeft
\newcommand{\DSource}[3]{\BeginSource{AB01-PDF#1}{#2}%
 \setlength{\FullSourceWidth}{\textwidth}\setlength{\GutterOffset}{0pt}%
 \ifodd#3\relax\NumbersLefttrue\else\NumbersLeftfalse\fi}
\newsavebox{\DLBox}
\newcommand{\DLine}[4]{%
 \par\noindent\hypertarget{#1}{}%
 \ifx\relax#2\relax\else
   \ifNumbersLeft
    \rlap{\kern-\GutterOffset\llap{\fontsize{7}{8}\selectfont #2\hspace{3mm}}}%
   \else
    \rlap{\kern\dimexpr\FullSourceWidth-\GutterOffset\relax\hspace{3mm}{\fontsize{7}{8}\selectfont #2}}%
   \fi
 \fi
 \ifx\relax#3\relax\else
   \ifNumbersLeft
    \rlap{\kern\dimexpr\FullSourceWidth-\GutterOffset\relax\hspace{3mm}{\fontsize{7}{8}\selectfont #3}}%
   \else
    \rlap{\kern-\GutterOffset\llap{\fontsize{7}{8}\selectfont #3\hspace{3mm}}}%
   \fi
 \fi
 \sbox{\DLBox}{\strut #4}%
 \ifdim\wd\DLBox>\linewidth\typeout{DIPLOMATICWIDTH|#1|\the\wd\DLBox|\the\linewidth}\fi
 \usebox{\DLBox}\strut\par}
\begin{document}
\thispagestyle{empty}
\begin{center}\Large AL-BATTĀNĪ\par\vspace{4mm}\LARGE OPUS ASTRONOMICUM\par
\vspace{5mm}\large Caroli Alphonsi Nallino pars secunda: frontispicium, praefatio,\\conspectus librorum, addenda et emendanda\par\vspace{8mm}
\Large S05 --- editio diplomatica\par\vspace{4mm}\large Paginae impressae [I]--VI\end{center}
\vspace{12mm}\noindent Singulae lineae paginarum impressarum singulis lineis huius editionis respondent, cum
divisionibus vocabulorum, signis et generibus litterarum (rectis, inclinatis, distantibus). Textus ex imaginibus paginarum transcriptus est;
litterae Graecae et Arabicae, numeri et signa in imaginibus amplificatis lecta sunt.
Haec transcriptio a nullo homine recognita est.\par
\DSource{0418}{[I]}{1}
\par\vspace*{\dimexpr 87.05mm-\baselineskip\relax}\prevdepth=0pt
\DLine{AB01-PDF0418-body-L001}{}{}{\CenterLine{\LFit{250.43pt}{14.01pt}{AL-BATTĀNĪ \textsc{sive} ALBATENII}}}
\par\vspace{7.09mm}
\DLine{AB01-PDF0418-body-L002}{}{}{\CenterLine{\LFit{336.01pt}{25.16pt}{OPUS ASTRONOMICUM}}}
\par\vspace{8.19mm}
\par\noindent\makebox[\linewidth][c]{\rule[.5ex]{32.43mm}{0.5pt}}\par
\DSource{0419}{[II]}{2}
\null
\DSource{0420}{[III]}{3}
\par\vspace*{\dimexpr 0.49mm-\baselineskip\relax}\prevdepth=0pt
\DLine{AB01-PDF0420-body-L001}{}{}{\CenterLine{\LFit{176.90pt}{15.82pt}{PUBBLICAZIONI}}}
\par\vspace{4.05mm}
\DLine{AB01-PDF0420-body-L002}{}{}{\CenterLine{\LFit{425.22pt}{14.77pt}{DEL REALE OSSERVATORIO DI BRERA IN MILANO.}}}
\par\vspace{3.52mm}
\DLine{AB01-PDF0420-body-L003}{}{}{\CenterLine{\LFit{128.23pt}{12.66pt}{\textbf{N. XL. Parte II.}}}}
\par\vspace{1.92mm}
\par\noindent\makebox[\linewidth][c]{\rule[.5ex]{118.54mm}{0.5pt}}\par
\par\vspace{5.92mm}
\DLine{AB01-PDF0420-body-L004}{}{}{\CenterLine{\LFit{295.78pt}{24.11pt}{AL-BATTĀNĪ \textsc{sive} ALBATENII}}}
\par\vspace{8.15mm}
\DLine{AB01-PDF0420-body-L005}{}{}{\CenterLine{\LFit{529.64pt}{35.33pt}{OPUS ASTRONOMICUM.}}}
\par\vspace{10.16mm}
\DLine{AB01-PDF0420-body-L006}{}{}{\CenterLine{\LFit{407.13pt}{9.94pt}{\textbf{AD FIDEM CODICIS ESCURIALENSIS ARABICE EDITUM}}}}
\par\vspace{5.85mm}
\DLine{AB01-PDF0420-body-L007}{}{}{\CenterLine{\LFit{261.58pt}{8.14pt}{LATINE VERSUM, ADNOTATIONIBUS INSTRUCTUM}}}
\par\vspace{5.04mm}
\DLine{AB01-PDF0420-body-L008}{}{}{\CenterLine{\LFit{6.63pt}{6.03pt}{\textbf{A}}}}
\par\vspace{3.72mm}
\DLine{AB01-PDF0420-body-L009}{}{}{\CenterLine{\LFit{348.97pt}{12.81pt}{\Name{CAROLO ALPHONSO NALLINO.}}}}
\par\vspace{19.74mm}
\par\noindent\makebox[\linewidth][c]{\rule[.5ex]{22.70mm}{1.1pt}}\par
\par\vspace{2.94mm}
\DLine{AB01-PDF0420-body-L010}{}{}{\CenterLine{\LFit{135.01pt}{11.75pt}{PARS SECUNDA}}}
\par\vspace{2.11mm}
\DLine{AB01-PDF0420-body-L011}{}{}{\CenterLine{\LFit{319.89pt}{10.25pt}{\textbf{VERSIO TABULARUM OMNIUM CUM ANIMADVERSIONIBUS,}}}}
\par\vspace{0.13mm}
\DLine{AB01-PDF0420-body-L012}{}{}{\CenterLine{\LFit{126.87pt}{10.25pt}{\textbf{GLOSSARIO, INDICIBUS.}}}}
\par\vspace{1.48mm}
\par\noindent\makebox[\linewidth][c]{\rule[.5ex]{22.54mm}{0.6pt}}\par
\par\vspace{28.71mm}
\DLine{AB01-PDF0420-body-L013}{}{}{\CenterLine{\LFit{177.65pt}{11.75pt}{\textbf{MEDIOLANI INSUBRUM}}}}
\par\vspace{1.06mm}
\DLine{AB01-PDF0420-body-L014}{}{}{\CenterLine{\LFit{268.66pt}{9.79pt}{P\textsc{rostat apud} \Name{ULRICHUM HOEPLIUM}}}}
\par\vspace{-0.95mm}
\DLine{AB01-PDF0420-body-L015}{}{}{\CenterLine{\LFit{79.11pt}{4.52pt}{\textbf{REGIUM BIBLIOPOLAM}}}}
\par\vspace{-0.57mm}
\DLine{AB01-PDF0420-body-L016}{}{}{\CenterLine{\LFit{193.17pt}{7.99pt}{\textbf{\textsc{in xysto christophoriano, nn.}} 58-63.}}}
\par\vspace{-0.49mm}
\par\noindent\makebox[\linewidth][c]{\rule[.5ex]{3.56mm}{0.95pt}}\par
\par\vspace{-0.70mm}
\DLine{AB01-PDF0420-body-L017}{}{}{\CenterLine{\LFit{28.78pt}{9.19pt}{1907.}}}
\DSource{0421}{[IV]}{4}
\par\vspace*{\dimexpr 216.35mm-\baselineskip\relax}\prevdepth=0pt
\par\noindent\makebox[\linewidth][c]{\rule[.5ex]{187.99mm}{0.6pt}}\par
\par\vspace{1.55mm}
\DLine{AB01-PDF0421-body-L001}{}{}{\CenterLine{\LFit{256.08pt}{9.34pt}{R\textsc{omae} MCMVII. — Excudebat C. D\textsc{e} L\textsc{uigi}.}}}
\DSource{0422}{[V]}{5}
\par\vspace*{\dimexpr -0.57mm-\baselineskip\relax}\prevdepth=0pt
\par\noindent\hypertarget{AB01-PDF0422-body-F01}{}\CenterLine{% Nallino, Part II p. V: the head ornament of the Praefatio, a heavy rule over a border of 33 equal sorts
% (a beaded bar, a lyre of two curled scrolls round a stem, side leaves and a hanging leaf); the sort
% averaged from its 33 printed copies and traced, set at the pitch fitted to them; the band drawn
% straight; enlarged in the ratio of the measures (_intake/latin/tools/fig422).
\begin{tikzpicture}[x=1.2557pt,y=-1.2557pt,line cap=butt,every node/.style={font=\fontsize{9}{9}\bfseries\itshape\selectfont,inner sep=0pt}]
\useasboundingbox (0,0) rectangle (426.48,24.97);
\fill (4.80,0.00) rectangle (425.98,2.16);  % the rule
\begin{scope}  % the border: the sort drawn 33 times, 12.8 pt apart
\clip (3.39,2.16) rectangle (425.98,24.97);
\foreach \k in {0,...,32} {\fill[even odd rule,shift={(\k*12.800,0)}]
  (9.87,22.30) -- (9.74,22.32) -- (9.56,22.29) -- (9.42,22.22) -- (9.31,22.12) -- (8.94,21.58) -- (8.68,20.98) -- (8.62,20.91) --
  (8.46,20.88) -- (8.11,21.00) -- (7.97,21.01) -- (7.75,20.93) -- (7.61,20.75) -- (7.49,20.45) -- (7.43,19.90) -- (7.51,19.32) --
  (7.71,18.70) -- (7.76,18.47) -- (7.77,18.30) -- (7.71,18.17) -- (7.62,18.13) -- (7.50,18.13) -- (6.92,18.46) -- (6.71,18.51) --
  (6.50,18.52) -- (6.31,18.49) -- (6.21,18.44) -- (6.10,18.29) -- (6.05,18.15) -- (6.07,18.00) -- (6.18,17.66) -- (6.20,17.32) --
  (6.15,16.93) -- (6.06,16.67) -- (6.00,16.30) -- (5.99,16.06) -- (6.06,15.56) -- (6.18,15.24) -- (6.29,15.04) -- (6.49,14.80) --
  (6.85,14.55) -- (7.15,14.41) -- (7.55,14.29) -- (7.99,14.29) -- (8.47,14.37) -- (8.58,14.34) -- (8.60,14.24) -- (8.56,13.66) --
  (8.62,13.39) -- (8.77,13.11) -- (9.05,12.85) -- (9.17,12.71) -- (9.27,12.50) -- (9.32,12.07) -- (9.32,11.80) -- (9.28,11.58) --
  (9.19,11.46) -- (9.01,11.38) -- (8.75,11.38) -- (8.62,11.42) -- (8.45,11.55) -- (8.29,11.73) -- (7.68,12.69) -- (7.36,13.13) --
  (7.07,13.48) -- (6.11,14.53) -- (5.56,15.16) -- (5.28,15.62) -- (4.96,16.25) -- (4.88,16.55) -- (4.86,16.79) -- (4.86,17.28) --
  (4.89,17.54) -- (4.97,17.80) -- (5.09,18.08) -- (5.18,18.35) -- (5.17,18.50) -- (5.07,18.68) -- (4.95,18.75) -- (4.73,18.74) --
  (4.64,18.71) -- (4.39,18.52) -- (4.19,18.24) -- (3.83,17.54) -- (3.68,16.85) -- (3.65,16.27) -- (3.68,15.98) -- (3.88,15.28) --
  (4.25,14.50) -- (4.28,14.31) -- (4.25,14.19) -- (4.17,14.11) -- (3.84,14.00) -- (3.72,13.91) -- (3.62,13.73) -- (3.60,13.55) --
  (3.70,13.20) -- (3.87,12.82) -- (3.96,12.38) -- (4.08,12.05) -- (4.28,11.70) -- (4.47,11.48) -- (4.64,11.32) -- (4.81,11.21) --
  (4.92,11.16) -- (5.32,11.07) -- (5.65,11.06) -- (5.92,11.08) -- (6.32,11.20) -- (6.67,11.43) -- (6.80,11.48) -- (6.89,11.49) --
  (7.06,11.46) -- (7.17,11.39) -- (7.27,11.24) -- (7.30,11.10) -- (7.25,10.96) -- (7.10,10.82) -- (6.42,10.48) -- (6.11,10.30) --
  (5.62,9.92) -- (5.38,9.70) -- (5.01,9.29) -- (4.77,8.98) -- (4.58,8.67) -- (4.31,8.06) -- (4.17,7.66) -- (4.07,7.30) --
  (3.98,6.64) -- (4.02,5.94) -- (4.14,5.25) -- (4.32,4.60) -- (4.29,4.58) -- (4.23,4.60) -- (3.88,4.82) -- (3.63,4.87) --
  (3.33,4.78) -- (3.13,4.66) -- (2.80,4.54) -- (2.77,4.22) -- (2.77,2.95) -- (2.78,2.75) -- (2.83,2.70) -- (5.15,2.70) --
  (5.76,2.68) -- (6.99,2.70) -- (7.66,2.68) -- (9.16,2.72) -- (9.95,2.70) -- (10.18,2.72) -- (10.74,2.68) -- (11.58,2.69) --
  (11.99,2.71) -- (12.48,2.69) -- (12.76,2.71) -- (14.49,2.71) -- (15.01,2.69) -- (16.72,2.70) -- (16.77,2.71) -- (16.80,2.74) --
  (16.81,4.39) -- (16.80,4.70) -- (16.78,4.79) -- (16.72,4.83) -- (16.62,4.85) -- (16.35,4.86) -- (16.08,4.78) -- (15.85,4.64) --
  (15.65,4.55) -- (15.57,4.55) -- (15.52,4.59) -- (15.76,5.47) -- (15.85,6.17) -- (15.87,6.58) -- (15.84,7.07) -- (15.75,7.60) --
  (15.58,8.14) -- (15.41,8.53) -- (15.00,9.19) -- (14.48,9.75) -- (13.66,10.42) -- (13.36,10.59) -- (12.64,10.94) -- (12.56,11.02) --
  (12.52,11.15) -- (12.55,11.27) -- (12.62,11.37) -- (12.73,11.42) -- (12.86,11.44) -- (13.02,11.41) -- (13.45,11.17) -- (13.66,11.11) --
  (14.07,11.10) -- (14.47,11.17) -- (14.77,11.28) -- (15.13,11.51) -- (15.30,11.68) -- (15.52,12.01) -- (15.67,12.40) -- (16.02,13.65) --
  (15.99,13.82) -- (15.88,13.94) -- (15.73,14.03) -- (15.41,14.10) -- (15.35,14.16) -- (15.33,14.24) -- (15.36,14.44) -- (15.74,15.23) --
  (15.88,15.66) -- (15.96,16.21) -- (15.95,16.68) -- (15.86,17.24) -- (15.73,17.65) -- (15.52,18.10) -- (15.30,18.41) -- (15.09,18.64) --
  (14.91,18.76) -- (14.80,18.78) -- (14.64,18.74) -- (14.52,18.65) -- (14.45,18.53) -- (14.44,18.40) -- (14.52,18.16) -- (14.71,17.79) --
  (14.80,17.52) -- (14.83,16.89) -- (14.81,16.68) -- (14.74,16.36) -- (14.43,15.73) -- (14.21,15.40) -- (12.68,13.61) -- (12.34,13.15) --
  (11.31,11.56) -- (11.17,11.46) -- (10.98,11.40) -- (10.71,11.42) -- (10.54,11.49) -- (10.45,11.62) -- (10.40,11.84) -- (10.42,12.23) --
  (10.47,12.50) -- (10.52,12.65) -- (10.62,12.83) -- (10.97,13.26) -- (11.06,13.62) -- (11.05,13.88) -- (10.92,14.33) -- (10.95,14.43) --
  (11.00,14.46) -- (11.10,14.46) -- (11.51,14.36) -- (11.91,14.32) -- (12.48,14.41) -- (12.89,14.61) -- (13.14,14.78) -- (13.39,15.07) --
  (13.50,15.27) -- (13.57,15.46) -- (13.61,15.74) -- (13.62,16.13) -- (13.42,17.15) -- (13.41,17.39) -- (13.46,17.65) -- (13.62,18.12) --
  (13.63,18.25) -- (13.59,18.38) -- (13.49,18.49) -- (13.39,18.55) -- (13.27,18.57) -- (12.98,18.54) -- (12.75,18.49) -- (12.56,18.41) --
  (12.12,18.18) -- (12.00,18.17) -- (11.92,18.21) -- (11.87,18.34) -- (11.87,18.54) -- (12.09,19.42) -- (12.15,19.82) -- (12.08,20.43) --
  (11.95,20.79) -- (11.84,20.93) -- (11.66,21.02) -- (11.57,21.05) -- (11.42,21.02) -- (11.06,20.84) -- (10.98,20.84) -- (10.89,20.88) --
  (10.81,20.98) -- (10.59,21.54) -- (10.33,21.96) -- (10.12,22.20) -- cycle
  (8.02,9.93) -- (8.19,9.85) -- (8.31,9.72) -- (8.79,8.71) -- (8.99,8.19) -- (9.26,6.88) -- (9.28,6.48) -- (9.27,6.30) --
  (9.23,6.24) -- (9.16,6.24) -- (8.56,6.72) -- (8.35,6.82) -- (8.02,6.94) -- (7.71,6.97) -- (7.33,6.95) -- (6.88,6.75) --
  (6.62,6.50) -- (6.52,6.36) -- (6.46,6.23) -- (6.39,5.96) -- (6.37,5.79) -- (6.40,5.54) -- (6.44,5.40) -- (6.56,5.18) --
  (6.67,5.04) -- (6.79,4.94) -- (6.99,4.83) -- (7.15,4.79) -- (7.37,4.77) -- (7.57,4.79) -- (7.82,4.85) -- (7.87,4.83) --
  (7.88,4.78) -- (7.85,4.70) -- (7.79,4.64) -- (7.61,4.50) -- (7.47,4.44) -- (7.11,4.42) -- (6.70,4.42) -- (6.54,4.45) --
  (6.23,4.56) -- (5.94,4.73) -- (5.61,5.03) -- (5.30,5.43) -- (5.11,5.92) -- (5.04,6.21) -- (5.05,6.89) -- (5.14,7.37) --
  (5.38,7.96) -- (5.56,8.25) -- (5.92,8.67) -- (6.38,9.15) -- (6.53,9.26) -- (7.29,9.74) -- (7.52,9.85) -- (7.75,9.92) --
  (7.88,9.94) -- cycle
  (12.17,9.93) -- (12.37,9.87) -- (12.58,9.76) -- (13.22,9.37) -- (13.58,9.09) -- (13.84,8.82) -- (14.21,8.40) -- (14.51,7.91) --
  (14.64,7.54) -- (14.77,7.03) -- (14.78,6.54) -- (14.73,5.97) -- (14.55,5.49) -- (14.40,5.21) -- (14.11,4.90) -- (13.76,4.65) --
  (13.46,4.50) -- (13.10,4.42) -- (12.69,4.43) -- (12.45,4.49) -- (12.30,4.56) -- (12.22,4.67) -- (12.27,4.71) -- (12.74,4.81) --
  (12.87,4.87) -- (13.16,5.08) -- (13.27,5.21) -- (13.36,5.39) -- (13.40,5.52) -- (13.42,5.73) -- (13.37,6.11) -- (13.23,6.41) --
  (13.07,6.60) -- (12.86,6.75) -- (12.61,6.90) -- (12.37,6.95) -- (12.08,6.97) -- (11.79,6.95) -- (11.60,6.90) -- (11.08,6.63) --
  (10.77,6.42) -- (10.69,6.43) -- (10.65,6.53) -- (10.63,6.96) -- (10.69,7.42) -- (10.77,7.79) -- (11.05,8.67) -- (11.30,9.31) --
  (11.44,9.59) -- (11.67,9.88) -- (11.88,9.95) -- cycle
  (11.03,3.93) -- (11.15,3.86) -- (11.31,3.72) -- (11.38,3.63) -- (11.37,3.56) -- (11.29,3.52) -- (11.03,3.48) -- (10.84,3.48) --
  (10.75,3.51) -- (10.70,3.57) -- (10.69,3.66) -- (10.70,3.79) -- (10.75,3.88) -- (10.86,3.95) -- cycle
  (9.00,3.92) -- (9.09,3.83) -- (9.15,3.73) -- (9.17,3.62) -- (9.14,3.54) -- (9.07,3.50) -- (8.95,3.49) -- (8.75,3.50) --
  (8.65,3.54) -- (8.62,3.64) -- (8.70,3.79) -- (8.88,3.92) -- (8.94,3.94) -- cycle
  (16.78,17.82) -- (16.69,17.67) -- (16.58,17.31) -- (16.48,16.83) -- (16.46,16.28) -- (16.50,15.93) -- (16.56,15.62) -- (16.71,15.21) --
  (16.78,15.09) -- (16.81,15.32) -- (16.81,16.92) -- (16.81,17.64) -- cycle
  (16.77,14.03) -- (16.60,13.98) -- (16.51,13.90) -- (16.45,13.78) -- (16.41,13.53) -- (16.51,13.21) -- (16.61,13.01) -- (16.70,12.75) --
  (16.79,12.29) -- (16.81,13.73) -- (16.80,13.96) -- cycle
  (2.80,17.93) -- (2.77,17.10) -- (2.77,15.16) -- (2.80,14.94) -- (2.90,15.13) -- (3.06,15.62) -- (3.15,16.15) -- (3.16,16.51) --
  (3.04,17.26) -- (2.91,17.70) -- cycle
  (4.79,4.04) -- (5.19,3.75) -- (5.26,3.67) -- (5.25,3.61) -- (5.18,3.58) -- (5.09,3.57) -- (4.88,3.65) -- (4.73,3.98) --
  (4.73,4.04) -- cycle
  (2.80,8.04) -- (2.77,7.24) -- (2.77,5.47) -- (2.79,4.83) -- (2.96,5.55) -- (3.04,6.10) -- (3.06,6.53) -- (3.02,7.07) --
  (2.94,7.56) -- cycle
  (10.50,10.12) -- (10.54,10.04) -- (10.54,9.93) -- (10.48,9.79) -- (10.42,9.75) -- (10.35,9.83) -- (10.34,10.00) -- (10.41,10.12) --
  (10.46,10.14) -- cycle
  (2.95,14.01) -- (2.81,14.03) -- (2.78,13.94) -- (2.77,12.78) -- (2.79,12.24) -- (3.01,12.86) -- (3.20,13.63) -- (3.19,13.74) --
  (3.15,13.84) -- cycle
  (15.06,3.88) -- (15.06,3.76) -- (14.96,3.65) -- (14.78,3.62) -- (14.66,3.67) -- (14.69,3.75) -- (14.84,3.85) -- (14.98,3.90) -- cycle;}
\end{scope}
\end{tikzpicture}
}\par
\par\vspace{38.31mm}
\DLine{AB01-PDF0422-body-L001}{}{}{\CenterLine{\LFit{183.67pt}{19.59pt}{\Name{PRAEFATIO}}}}
\par\vspace{7.21mm}
\par\noindent\makebox[\linewidth][c]{\rule[.5ex]{28.17mm}{0.48pt}}\par
\par\fontsize{9.4}{10.82}\selectfont
\par\vspace{15.68mm}
\DLine{AB01-PDF0422-body-L002}{}{}{\textbf{§ 1, Quali ratione tabulas emendaverim. — § 2, Nova de versione Hispanica eiusque codice Parisiensi.}}
\DLine{AB01-PDF0422-body-L003}{}{}{\mbox{}\rlap{\hspace*{10.26mm}\smash{\rlap{\textbf{— § 3, Fontes et ordo tabularum.}}}}}
\par\fontsize{11.1}{14.81}\selectfont
\par\vspace{4.60mm}
\hypertarget{AB01-PDF0422-P01}{}
\DLine{AB01-PDF0422-body-L004}{}{}{\Indent \textbf{§ 1.} — Monui interdum in notis meis numeros codicis Escurialensis crebris scatere men-}
\par\vspace{0.01mm}
\DLine{AB01-PDF0422-body-L005}{}{}{dis, ita ut nulla tabula in codice occurrat erroribus destituta. Ex 180 sinibus, ut exemplum}
\par\vspace{0.08mm}
\DLine{AB01-PDF0422-body-L006}{}{}{praebeam, 74 in minutis aut secundis mendosi sunt. Pigebat tabulas vitiatas, quales in codice}
\par\vspace{-0.03mm}
\DLine{AB01-PDF0422-body-L007}{}{}{exstabant, proferre; taedebat omnes lectiones falsas recensere quae nullius utilitatis erant;}
\DLine{AB01-PDF0422-body-L008}{}{}{visum est igitur tabulas pro viribus meis expurgatas edere, iis tantum lectionibus codicis}
\DLine{AB01-PDF0422-body-L009}{}{}{indicatis quarum emendatio minus certa aut perspicua.}
\par\vspace{0.04mm}
\hypertarget{AB01-PDF0422-P02}{}
\DLine{AB01-PDF0422-body-L010}{}{}{\Indent Astronomi Arabum numeros tabularum non figuris quas nunc Arabicas vocamus scribunt,}
\DLine{AB01-PDF0422-body-L011}{}{}{sed litteris alphabeti quibus valores numerales tribuunt. E permutationibus similium litterarum,}
\DLine{AB01-PDF0422-body-L012}{}{}{festinanter aut negligenter a librariis exaratarum, facile corruptelae oriuntur. Litterae cre-}
\par\vspace{0.07mm}
\DLine{AB01-PDF0422-body-L013}{}{}{brius inter se commutatae hoc nostro codice, characteribus Arabico–occidentalibus seu maghre-}
\par\vspace{0.16mm}
\DLine{AB01-PDF0422-body-L014}{}{}{binis scripto, haec sunt:}
\par\vspace{2.21mm}
\DLine{AB01-PDF0422-body-L015}{}{}{\mbox{}\rlap{\hspace*{37.51mm}\smash{\llap{\ArabicRun{ب}}}}\rlap{\hspace*{42.30mm}\smash{\rlap{(\ArabicRun{ٮ}) 2 cum \ArabicRun{ى} 10 aut \ArabicRun{د} 4.}}}\rlap{\hspace*{112.01mm}\smash{\llap{\ArabicRun{ج}}}}\rlap{\hspace*{116.91mm}\smash{\rlap{3 cum \ArabicRun{ح} 8.}}}}
\DLine{AB01-PDF0422-body-L016}{}{}{\mbox{}\rlap{\hspace*{37.51mm}\smash{\llap{\ArabicRun{د}}}}\rlap{\hspace*{42.30mm}\smash{\rlap{4 cum \ArabicRun{ه} 5 aut \ArabicRun{و} 6.}}}\rlap{\hspace*{112.01mm}\smash{\llap{\ArabicRun{ه}}}}\rlap{\hspace*{116.91mm}\smash{\rlap{5 cum \ArabicRun{ط} 9.}}}}
\DLine{AB01-PDF0422-body-L017}{}{}{\mbox{}\rlap{\hspace*{37.51mm}\smash{\llap{\ArabicRun{ز}}}}\rlap{\hspace*{42.30mm}\smash{\rlap{(\ArabicRun{ر}) 7 cum \ArabicRun{و} 6 aut \ArabicRun{ن} 50.}}}\rlap{\hspace*{112.01mm}\smash{\llap{\ArabicRun{ط}}}}\rlap{\hspace*{116.91mm}\smash{\rlap{9 cum \ArabicRun{ك} 20.}}}}
\par\vspace{0.10mm}
\DLine{AB01-PDF0422-body-L018}{}{}{\mbox{}\rlap{\hspace*{37.51mm}\smash{\llap{\ArabicRun{ي‍}}}}\rlap{\hspace*{42.30mm}\smash{\rlap{(\ArabicRun{ٮ‍}) 10 cum \ArabicRun{ل‍} 30 aut \ArabicRun{ن‍} 50.}}}\rlap{\hspace*{112.01mm}\smash{\llap{\ArabicRun{ك‍}}}}\rlap{\hspace*{116.91mm}\smash{\rlap{20 cum \ArabicRun{ل‍} 30.}}}}
\DLine{AB01-PDF0422-body-L019}{}{}{\mbox{}\rlap{\hspace*{37.51mm}\smash{\llap{\ArabicRun{ن‍}}}}\rlap{\hspace*{42.30mm}\smash{\rlap{50 cum \ArabicRun{ڢ‍} 80 aut \ArabicRun{ف‍} 100.}}}\rlap{\hspace*{112.01mm}\smash{\llap{\ArabicRun{‍ع‍}}}}\rlap{\hspace*{116.91mm}\smash{\rlap{70 cum \ArabicRun{‍ڡ‍} (\ArabicRun{‍ڢ‍}) 80.}}}}
\par\vspace{0.17mm}
\DLine{AB01-PDF0422-body-L020}{}{}{\mbox{}\rlap{\hspace*{37.51mm}\smash{\llap{\ArabicRun{ص}}}}\rlap{\hspace*{42.30mm}\smash{\rlap{60 cum \ArabicRun{ض} 90.}}}}
\par\vspace{2.18mm}
\hypertarget{AB01-PDF0422-P03}{}
\DLine{AB01-PDF0422-body-L021}{}{}{\Indent Rarius mendae ex eo proveniunt quod alter est valor numeralis nonnullarum litterarum}
\DLine{AB01-PDF0422-body-L022}{}{}{in scriptura maghrebina codicis, alter in scriptura Arabico–orientali seu naskhī qua al-Battānī}
\DLine{AB01-PDF0422-body-L023}{}{}{procul dubio est usus. Qui archetypum maghrebinum scripsit, unde codex Escurialensis codi-}
\par\vspace{-0.10mm}
\DLine{AB01-PDF0422-body-L024}{}{}{cesque a Platone Tiburtino et ab interprete Hispano adhibiti fluxerunt, interdum (raro tamen)}
\DLine{AB01-PDF0422-body-L025}{}{}{opportunas mutationes negligens numeros vitiavit. Nam \ArabicRun{س} 60, \ArabicRun{ص} 90, \ArabicRun{ف} 80, \ArabicRun{ض} 800, \ArabicRun{ظ} 900}
\DLine{AB01-PDF0422-body-L026}{}{}{Albatenii significant 300, 60, 100, 90 et 800 in codice Escurialensi. Confusionis exemplum}
\par\vspace{0.19mm}
\DLine{AB01-PDF0422-body-L027}{}{}{memoravi t. III, p. 230, adǹ. 8.}
\DSource{0423}{VI}{6}
\par\fontsize{11.1}{14.92}\selectfont
\par\vspace*{\dimexpr 0.09mm-\baselineskip\relax}\prevdepth=0pt
\hypertarget{AB01-PDF0423-P01}{}
\DLine{AB01-PDF0423-body-L001}{}{}{\Indent Rarius etiam errores ex exemplari Arabico–orientali recepti videntur, ut commutationes}
\par\vspace{0.02mm}
\DLine{AB01-PDF0423-body-L002}{}{}{\ArabicRun{ب} 2 cum \ArabicRun{ج} 3, vel \ArabicRun{ج} 3 cum \ArabicRun{د} 4; Arabes enim orientales in libris mathematicis numerum 3}
\par\vspace{0.04mm}
\DLine{AB01-PDF0423-body-L003}{}{}{scribunt omissis puncto et cauda (1), scilicet \ArabicRun{ح‍} (sed 8 semper \ArabicRun{ح}). Nostrates, qui hanc rem non}
\DLine{AB01-PDF0423-body-L004}{}{}{animadverterunt, in edendis aut vertendis libris Arabum orientalium perfrequenter errave-}
\DLine{AB01-PDF0423-body-L005}{}{}{runt ac 8 pro 3 posuerunt.}
\par\vspace{0.05mm}
\hypertarget{AB01-PDF0423-P02}{}
\DLine{AB01-PDF0423-body-L006}{}{}{\Indent Septies pro \AbjadZero{} 0 negationem Arabicam \ArabicRun{لا} \textit{lā} scriptam invenimus, quae, iuxta litterarum}
\DLine{AB01-PDF0423-body-L007}{}{}{valorem, 31 significat; idque p. 53 (in minutis latit. urbis as-Siraǵān), 57 (in secundis decli-}
\par\vspace{-0.04mm}
\DLine{AB01-PDF0423-body-L008}{}{}{nationis pro 1°), 84 (in secundis motus latitudinis, ult. lin.), 93 (in secundis tertii termini pro}
\DLine{AB01-PDF0423-body-L009}{}{}{20° et 32°), 160 (in gradibus latit. octavae stellae Virginis), et 185 (in gradibus cum quibus}
\DLine{AB01-PDF0423-body-L010}{}{}{nr. 63 oritur). Eadem res interdum in exemplari occurrebat quod interpres Hispanus adhibuit,}
\DLine{AB01-PDF0423-body-L011}{}{}{ut t. II, p. 215, adn. 66, p. 217, adn. 261, et p. 271 monui. Cfr. t. II, p. 266, l. 5–9, ubi talis}
\DLine{AB01-PDF0423-body-L012}{}{}{erroris exemplum ex \textit{Tabulis Alphonsinis} praebui.}
\hypertarget{AB01-PDF0423-P03}{}
\DLine{AB01-PDF0423-body-L013}{}{}{\Indent Corruptelas sanandas quod attinet, triplex est tabularum genus. \textit{Primum} tabularum qua-}
\DLine{AB01-PDF0423-body-L014}{}{}{rum numeri ab alteris deducuntur per additiones aut subductiones datae quantitatis; tunc er-}
\DLine{AB01-PDF0423-body-L015}{}{}{ror numerum in codice afficiens si numeris sequentibus non propagatur, manifeste librariis est}
\par\vspace{-0.05mm}
\DLine{AB01-PDF0423-body-L016}{}{}{tribuendus ac facillime emendandus. \textit{Alterum} habetur genus cum inter numeros tabulae ratio}
\DLine{AB01-PDF0423-body-L017}{}{}{quaedam sit, tamen non adeo certa ut possint discrepantiae unius vel duarum secundarum}
\DLine{AB01-PDF0423-body-L018}{}{}{(si secundae exhibeantur), unius vel duorum minutorum (si haec tantum occurrant cum gra-}
\DLine{AB01-PDF0423-body-L019}{}{}{dibus) ad mentem auctoris tutissime restitui; aut etiam cum numeri nullam rationem mani-}
\DLine{AB01-PDF0423-body-L020}{}{}{festam inter se habeant ac singillatim sint supputatione trigonometrica deducendi, in qua di-}
\DLine{AB01-PDF0423-body-L021}{}{}{screpantiae unius vel duarum unitatum in secundis aut in minutis possunt ab ipsis auctoris}
\par\vspace{0.02mm}
\DLine{AB01-PDF0423-body-L022}{}{}{calculis provenire. \textit{Tertium} genus tunc habetur cum numeri inter se nullam rationem habeant}
\DLine{AB01-PDF0423-body-L023}{}{}{nec calculis deduci possint (ut in tabulis geographicis et catalogo fixarum), qua re tabulae}
\par\vspace{0.03mm}
\DLine{AB01-PDF0423-body-L024}{}{}{nonnisi proxime emendabiles sint per collationem fontium e quibus manaverint.}
\par\vspace{0.01mm}
\hypertarget{AB01-PDF0423-P04}{}
\DLine{AB01-PDF0423-body-L025}{}{}{\Indent Errores \textit{primi} generis tacite emendavi. Errores \textit{alterius} certissime sanavi quotiescumque}
\par\vspace{-0.01mm}
\DLine{AB01-PDF0423-body-L026}{}{}{Ptolemaei aut Theonis comparatio succurrebat; his subsidiis deficientibus, ut in declinationibus}
\par\vspace{-0.11mm}
\DLine{AB01-PDF0423-body-L027}{}{}{Solis et ascensionibus signorum, oportuit numeros logarithmis supputare, ad quorum fidem}
\DLine{AB01-PDF0423-body-L028}{}{}{discrepantiae maiores una vel duabus vel tribus etiam secundis emendare licebat. Numeros}
\DLine{AB01-PDF0423-body-L029}{}{}{falsos tabularum \textit{tertii} generis quantum potui ad mentem auctoris restituendos curavi quoties-}
\par\vspace{0.07mm}
\DLine{AB01-PDF0423-body-L030}{}{}{cumque restitutiones parvis ac facillimis litterarum permutationibus obtinebantur; semper au-}
\par\vspace{0.02mm}
\DLine{AB01-PDF0423-body-L031}{}{}{tem in notis numeros codicis exposui. — Uberiora de correctionibus adhibitis in commentario}
\DLine{AB01-PDF0423-body-L032}{}{}{in singulas tabulas narravi.}
\par\vspace{0.09mm}
\hypertarget{AB01-PDF0423-P05}{}
\DLine{AB01-PDF0423-body-L033}{}{}{\Indent In hoc taediosissimo labore tot milia numerorum probandi, summa qua est benevolentia}
\par\vspace{-0.07mm}
\DLine{AB01-PDF0423-body-L034}{}{}{socius mihi \Name{Ioh. Schiaparelli} fuit, qui tabulas a me primo emendatas iterum diligentis-}
\par\vspace{-0.02mm}
\DLine{AB01-PDF0423-body-L035}{}{}{sime examinavit. Solus Schiaparelli tabulas expurgavit quibus declinationes Solis (p. 57-58),}
\par\vspace{-0.03mm}
\DLine{AB01-PDF0423-body-L036}{}{}{longitudines dierum (p. 59), status fixarum anno 1211 Dhū-l-qarnayn (p. 178-186), altitudines}
\DLine{AB01-PDF0423-body-L037}{}{}{et directiones horarum (p. 188) exhibentur.}
\par\vspace{6pt}\nointerlineskip\hbox to\textwidth{\hss\vrule width 115.43pt height.25pt depth0pt\hss}\prevdepth=-1000pt \vspace{5pt}
\noindent\begin{minipage}[t]{.48\textwidth}\vspace{0pt}\fontsize{9}{11.5}\selectfont\setlength{\GutterOffset}{0pt}\setlength{\FullSourceWidth}{\linewidth}
\par\fontsize{9}{11.75}\selectfont
\par\vspace{-0.94mm}
\hypertarget{AB01-PDF0423-N01}{}
\DLine{AB01-PDF0423-notes_left-L001}{}{}{\Indent (1) Ita apud Ḥabash, Ibn Yūnus, Kūshiyār ibn}
\par\vspace{-0.02mm}
\DLine{AB01-PDF0423-notes_left-L002}{}{}{Labbān, Abū ’l-fidā’ etc. Cfr. \Name{Muḥammad Sibṭ}}
\DLine{AB01-PDF0423-notes_left-L003}{}{}{\Name{al-Māridīnī} apud \Name{Woepcke}, \textit{Sur l’introdu-}}
\par\vspace{0.03mm}
\DLine{AB01-PDF0423-notes_left-L004}{}{}{\textit{ction de l’arithmétique indienne en Occident}, Ro-}
\DLine{AB01-PDF0423-notes_left-L005}{}{}{me 1859, p. 69: in litteris alphabeti quae ad numeros}
\DLine{AB01-PDF0423-notes_left-L006}{}{}{fractionum sexagesimarum designandos inserviunt,}
\DLine{AB01-PDF0423-notes_left-L007}{}{}{« il ne peut y avoir de l’incertitude que relative-}
\end{minipage}
\hfill\vrule width.3pt\hfill\begin{minipage}[t]{.48\textwidth}\vspace{0pt}\fontsize{9}{11.5}\selectfont\setlength{\GutterOffset}{0pt}\setlength{\FullSourceWidth}{\linewidth}
\par\fontsize{9}{11.75}\selectfont
\par\vspace{-0.92mm}
\DLine{AB01-PDF0423-notes_right-L001}{}{}{« ment à deux de ces lettres, qui sont le \textit{noûn} (qu’on}
\par\vspace{-0.02mm}
\DLine{AB01-PDF0423-notes_right-L002}{}{}{« pourrait confondre) avec le \textit{yâ}; et le \textit{djîm} (qu’on}
\DLine{AB01-PDF0423-notes_right-L003}{}{}{« pourrait confondre) avec le \textit{hâ}. On est convenu,}
\par\vspace{0.02mm}
\DLine{AB01-PDF0423-notes_right-L004}{}{}{« à cause de cela, de donner toujours le point au}
\DLine{AB01-PDF0423-notes_right-L005}{}{}{« \textit{noûn}, à l’exclusion des autres (lettres), et à cou-}
\par\vspace{0.03mm}
\DLine{AB01-PDF0423-notes_right-L006}{}{}{« per le \textit{djîm} de la manière suivante: \ArabicRun{ح‍} ».}
\end{minipage}
\end{document}
